% GENERATED FILE. Edit canonical lessons, then run npm run generate:books.
% canonical-source-hash: fnv1a64:ac1869638dcd80b0
% canonical-lessons: ML-C472-gatagatam, ML-C472-gatagatakkurukku, ML-C472-vegata, ML-C472-nirannuka, ML-C472-kulunnuka

\chapter{Traffic}
\label{ch:ml-traffic}
\hlchaptermodality{472}

\begin{chapteropening}
\textbf{By the end of this chapter:} \emph{I can say transport, a traffic jam, speed, to creep along and to jolt.}

Complete the last lesson of chapter 472: I can say transport, a traffic jam, speed, to creep along and to jolt.
\end{chapteropening}

\section[\texorpdfstring{\ml{ഗതാഗതം} (gatāgataṁ)}{ഗതാഗതം (gatāgataṁ)}]{\ml{ഗതാഗതം} (gatāgataṁ) — transport; traffic}

\label{lesson:ML-C472-gatagatam}



\begin{quote}
Before the new one: say the Malayalam for a horn, then the Malayalam for to get dented.
\end{quote}

\subsection*{You'll want to know: \ml{ഗതാഗതം}}
\textbf{\ml{ഗതാഗതം}} — \emph{gatāgataṁ} — \textquotedblleft{}transport; traffic\textquotedblright{}.

\begin{quote}
\ml{പുതിയ} \ml{പാലം} \ml{ഗതാഗതം} \ml{എളുപ്പമാക്കി}.
\end{quote}

\begin{quote}
\emph{putiya pālaṁ gatāgataṁ eḷuppamākki.}
\end{quote}

\begin{quote}
The new bridge made transport easier.
\end{quote}

A formal word for the movement of vehicles and people; it is common in the news.

\begin{cousinweb}
\textbf{\ml{ഗതാഗതം}} (\emph{gatāgataṁ}) is borrowed from Sanskrit \emph{gatāgata}, \textquotedblleft{}going and coming\textquotedblright{} (Monier-Williams), a compound of \emph{gata}, gone, and \emph{āgata}, come.
\end{cousinweb}

\begin{grammarlens}[title={What you've built}]
One more word for getting around.
\end{grammarlens}

\subsection*{Your turn}
\begin{itemize}
\raggedright
  \item \emph{Say it:} \emph{gatāgataṁ}
  \item \emph{Say it:} \emph{gatāgataṁ}, once more
  \item \emph{Say it:} \emph{putiya pālaṁ gatāgataṁ eḷuppamākki.} — The new bridge made transport easier.
  \item \emph{Recall:} say the Malayalam for a horn, then the Malayalam for to get dented, then say \emph{gatāgataṁ} again
\end{itemize}

\begin{tcolorbox}[breakable,colback=teal!4,colframe=teal!35!black,title={Before you move on}]
What does \ml{ഗതാഗതം} mean? (\textquotedblleft{}Transport; traffic\textquotedblright{}.) Say it once more.
\end{tcolorbox}
\section[\texorpdfstring{\ml{ഗതാഗതക്കുരുക്ക്} (gatāgatakkurukkŭ)}{ഗതാഗതക്കുരുക്ക് (gatāgatakkurukkŭ)}]{\ml{ഗതാഗതക്കുരുക്ക്} (gatāgatakkurukkŭ) — a traffic jam}

\label{lesson:ML-C472-gatagatakkurukku}



\begin{quote}
Before the new one: say the Malayalam for to get dented, then the Malayalam for transport.
\end{quote}

\subsection*{You'll want to know: \ml{ഗതാഗതക്കുരുക്ക്}}
\textbf{\ml{ഗതാഗതക്കുരുക്ക്}} — \emph{gatāgatakkurukkŭ} — \textquotedblleft{}a traffic jam\textquotedblright{}.

\begin{quote}
\ml{ഓഫീസിലേക്കുള്ള} \ml{വഴിയിൽ} \ml{വലിയ} \ml{ഗതാഗതക്കുരുക്ക്} \ml{ഉണ്ടായിരുന്നു}.
\end{quote}

\begin{quote}
\emph{ōphīsilēkkuḷḷa vaḻiyil valiya gatāgatakkurukkŭ uṇṭāyirunnu.}
\end{quote}

\begin{quote}
There was a big traffic jam on the way to the office.
\end{quote}

\ml{കുരുക്ക്} is a knot or a snare: traffic tied in a knot. In speech people often say \ml{ബ്ലോക്ക്}.

\begin{cousinweb}
\textbf{\ml{ഗതാഗതക്കുരുക്ക്}} (\emph{gatāgatakkurukkŭ}) is a Malayalam compound built from \ml{ഗതാഗതം}, from Sanskrit \emph{gatāgata} (Monier-Williams), and \ml{കുരുക്ക്}, a knot or noose, whose origin is not settled.
\end{cousinweb}

\begin{grammarlens}[title={What you've built}]
One more word for getting around.
\end{grammarlens}

\subsection*{Your turn}
\begin{itemize}
\raggedright
  \item \emph{Say it:} \emph{gatāgatakkurukkŭ}
  \item \emph{Say it:} \emph{gatāgatakkurukkŭ}, once more
  \item \emph{Say it:} \emph{ōphīsilēkkuḷḷa vaḻiyil valiya gatāgatakkurukkŭ uṇṭāyirunnu.} — There was a big traffic jam on the way to the office.
  \item \emph{Recall:} say the Malayalam for to get dented, then the Malayalam for transport, then say \emph{gatāgatakkurukkŭ} again
\end{itemize}

\begin{tcolorbox}[breakable,colback=teal!4,colframe=teal!35!black,title={Before you move on}]
What does \ml{ഗതാഗതക്കുരുക്ക്} mean? (\textquotedblleft{}A traffic jam\textquotedblright{}.) Say it once more.
\end{tcolorbox}
\section[\texorpdfstring{\ml{വേഗത} (vēgata)}{വേഗത (vēgata)}]{\ml{വേഗത} (vēgata) — speed}

\label{lesson:ML-C472-vegata}



\begin{quote}
Before the new one: say the Malayalam for transport, then the Malayalam for a traffic jam.
\end{quote}

\subsection*{You'll want to know: \ml{വേഗത}}
\textbf{\ml{വേഗത}} — \emph{vēgata} — \textquotedblleft{}speed\textquotedblright{}.

\begin{quote}
\ml{ഇവിടെ} \ml{വണ്ടികളുടെ} \ml{വേഗത} \ml{കുറവാണ്}.
\end{quote}

\begin{quote}
\emph{iviṭe vaṇṭikaḷuṭe vēgata kuṟavāṇŭ.}
\end{quote}

\begin{quote}
The speed of vehicles is low here.
\end{quote}

\ml{വേഗം} is quickness; \ml{വേഗത} is speed as a measure. Fast is \ml{വേഗത്തിൽ}.

\begin{cousinweb}
\textbf{\ml{വേഗത}} (\emph{vēgata}) is built from Sanskrit \emph{vega}, \textquotedblleft{}a rush, speed\textquotedblright{} (Monier-Williams), with the ending \emph{-tā}, which makes an abstract noun.
\end{cousinweb}

\begin{grammarlens}[title={What you've built}]
One more word for getting around.
\end{grammarlens}

\subsection*{Your turn}
\begin{itemize}
\raggedright
  \item \emph{Say it:} \emph{vēgata}
  \item \emph{Say it:} \emph{vēgata}, once more
  \item \emph{Say it:} \emph{iviṭe vaṇṭikaḷuṭe vēgata kuṟavāṇŭ.} — The speed of vehicles is low here.
  \item \emph{Recall:} say the Malayalam for transport, then the Malayalam for a traffic jam, then say \emph{vēgata} again
\end{itemize}

\begin{tcolorbox}[breakable,colback=teal!4,colframe=teal!35!black,title={Before you move on}]
What does \ml{വേഗത} mean? (\textquotedblleft{}Speed\textquotedblright{}.) Say it once more.
\end{tcolorbox}
\section[\texorpdfstring{\ml{നിരങ്ങുക} (niraṅṅuka)}{നിരങ്ങുക (niraṅṅuka)}]{\ml{നിരങ്ങുക} (niraṅṅuka) — to creep along; to crawl}

\label{lesson:ML-C472-nirannuka}



\begin{quote}
Before the new one: say the Malayalam for a traffic jam, then the Malayalam for speed.
\end{quote}

\subsection*{You'll want to know: \ml{നിരങ്ങുക}}
\textbf{\ml{നിരങ്ങുക}} — \emph{niraṅṅuka} — \textquotedblleft{}to creep along; to crawl\textquotedblright{}.

\begin{quote}
\ml{തിരക്കിൽ} \ml{ബസ്} \ml{പതുക്കെ} \ml{നിരങ്ങി} \ml{നീങ്ങി}.
\end{quote}

\begin{quote}
\emph{tirakkil basŭ patukke niraṅṅi nīṅṅi.}
\end{quote}

\begin{quote}
In the crowd the bus crept slowly forward.
\end{quote}

Past: \ml{നിരങ്ങി}, \emph{niraṅṅi}.

Often paired with \ml{നീങ്ങുക}, to move: \ml{നിരങ്ങി} \ml{നീങ്ങുക}, to move at a crawl.

\begin{cousinweb}
\textbf{\ml{നിരങ്ങുക}} (\emph{niraṅṅuka}) is native Dravidian, \textquotedblleft{}to drag the feet along the ground, to creep\textquotedblright{} (DEDR 3749).
\end{cousinweb}

\begin{grammarlens}[title={What you've built}]
One more word for getting around.
\end{grammarlens}

\subsection*{Your turn}
\begin{itemize}
\raggedright
  \item \emph{Say it:} \emph{niraṅṅuka}
  \item \emph{Say it:} \emph{niraṅṅuka}, once more
  \item \emph{Say it:} \emph{tirakkil basŭ patukke niraṅṅi nīṅṅi.} — In the crowd the bus crept slowly forward.
  \item \emph{Recall:} say the Malayalam for a traffic jam, then the Malayalam for speed, then say \emph{niraṅṅuka} again
\end{itemize}

\begin{tcolorbox}[breakable,colback=teal!4,colframe=teal!35!black,title={Before you move on}]
What does \ml{നിരങ്ങുക} mean? (\textquotedblleft{}To creep along; to crawl\textquotedblright{}.) Say it once more.
\end{tcolorbox}
\section[\texorpdfstring{\ml{കുലുങ്ങുക} (kuluṅṅuka)}{കുലുങ്ങുക (kuluṅṅuka)}]{\ml{കുലുങ്ങുക} (kuluṅṅuka) — to jolt; to shake}

\label{lesson:ML-C472-kulunnuka}



\begin{quote}
Before the new one: say the Malayalam for speed, then the Malayalam for to creep along.
\end{quote}

\subsection*{You'll want to know: \ml{കുലുങ്ങുക}}
\textbf{\ml{കുലുങ്ങുക}} — \emph{kuluṅṅuka} — \textquotedblleft{}to jolt; to shake\textquotedblright{}.

\begin{quote}
\ml{മോശം} \ml{വഴിയിൽ} \ml{ബസ്} \ml{വല്ലാതെ} \ml{കുലുങ്ങി}.
\end{quote}

\begin{quote}
\emph{mōśaṁ vaḻiyil basŭ vallāte kuluṅṅi.}
\end{quote}

\begin{quote}
On the bad road the bus jolted badly.
\end{quote}

Past: \ml{കുലുങ്ങി}, \emph{kuluṅṅi}.

The partner of \ml{കുലുക്കുക}, to shake something: the bus \ml{കുലുങ്ങുക} by itself.

\begin{cousinweb}
\textbf{\ml{കുലുങ്ങുക}} (\emph{kuluṅṅuka}) is native Dravidian, the intransitive partner of \ml{കുലുക്കുക}, to shake (DEDR 1806).
\end{cousinweb}

\begin{grammarlens}[title={What you've built}]
One more word for getting around.
\end{grammarlens}

\subsection*{Your turn}
\begin{itemize}
\raggedright
  \item \emph{Say it:} \emph{kuluṅṅuka}
  \item \emph{Say it:} \emph{kuluṅṅuka}, once more
  \item \emph{Say it:} \emph{mōśaṁ vaḻiyil basŭ vallāte kuluṅṅi.} — On the bad road the bus jolted badly.
  \item \emph{Recall:} say the Malayalam for speed, then the Malayalam for to creep along, then say \emph{kuluṅṅuka} again
\end{itemize}

\begin{tcolorbox}[breakable,colback=teal!4,colframe=teal!35!black,title={Before you move on}]
What does \ml{കുലുങ്ങുക} mean? (\textquotedblleft{}To jolt; to shake\textquotedblright{}.) Say it once more.
\end{tcolorbox}
